%%%%%%%%%%%%%%%%%%%%%%%%%%%%%%%%%%%%%%%%%%%%%%%%%%%%%%%%%%%%%%%%%%%%
%%%
%%%  Creates the correct "Table of Contents" for your proceedings
%%%
%%%%%%%%%%%%%%%%%%%%%%%%%%%%%%%%%%%%%%%%%%%%%%%%%%%%%%%%%%%%%%%%%%%%
%%%
%%% Usage:
%%%    Paste the following commands into your LaTeX file :
%%%
%%%        \input proctoc    %%% this file
%%%        \input <database>    %%% the list of articles
%%%        \tableofcontents      %%% generate the table of contents
%%%
%%%    and replace <database> by the name of a texfile,
%%%    containing for each article an entry of the format
%%%
%%%     \article[<optRemark>]{<paperId>}{noPages}{author}{title}
%%%
%%%%%%%%%%%%%%%%%%%%%%%%%%%%%%%%%%%%%%%%%%%%%%%%%%%%%%%%%%%%%%%%%%%
%%%
%%%  Macro for formatting each entry in the table of contents (toc)
%%%
%%%  usage:
%%%     \formatEntry[optionalRemark]{author(s)}{title}{pageNumber}
%%%
%%%  redefine this to suit your taste
%%%
%%%  Example:
%%%
%%%    \newcommand{\formatEntry}[4]{
%%%         \contentsline{section}{{\sc #2} {\rm #1} \\ {\rm #3 \dotfill}}{\rm #4}}

\newcommand{\formatEntry}[4]{
    \contentsline{section}{{\sc #2} {\rm #1} \\ {\rm #3 \dotfill}}{\rm #4}}

%%%%%%%%%%%%%%%%%%%%%%%%%%%%%%%%%%%%%%%%%%%%%%%%%%%%%%%%%%%%%%%%%%
%%%
%%%      Macros for calculating the number
%%%      of the starting page for each article
%%%      and for generating its entry in the
%%%      table of contents
%%%
%%%%%%%%%%%%%%%%%%%%%%%%%%%%%%%%%%%%%%%%%%%%%%%%%%%%%%%%%%%%%%%%%%

 \newcounter{startPage}
 \setcounter{startPage}{1}

 \newcommand{\article}[5][]{
     \newcounter{#2}
     \setcounter{#2}{\value{startPage}}
     \addtocontents{toc}
        {\formatEntry {#1}{#4}{ #5}{\thestartPage}}
     \addtocounter{startPage}{#3}
 }
